\chapter{Ուսանող 1 — Միկրոէլեկտրոնիկա}

\textbf{Անհատական թեմա.}\\
\emph{AUP ZU3 FPGA-ի հիման վրա սենսորային տվյալների գեներացման, հավաքագրման և
նախնական մշակման ապարատային համակարգի նախագծում.}

\textbf{Հիմնական պարտականությունները.}
\begin{enumerate}[leftmargin=1.6em]
  \item Կոնֆիգուրելի RTL ազդանշան-գեներատորների նախագծում։
  \item Սենսոր-էմուլյացիայի տրամաբանության իրականացում։
  \item LFSR/աղմուկի գեներացիայի իրականացում։
  \item Անոմալիաների/անսարքությունների ներարկման իրականացում։
  \item Տվյալների հավաքագրման (data acquisition) տրամաբանության մշակում։
  \item FIFO/BRAM բուֆերացման իրականացում։
  \item Կոնֆիգուրելի processing IP-ների մշակում։
  \item AXI-ի հիման վրա տվյալների հաղորդակցության իրականացում։
  \item AXI4-Lite կոնֆիգուրացիոն ռեգիստրների իրականացում։
  \item PL-ի ինտեգրում Zynq PS-ի հետ։
  \item Testbench-երի մշակում և RTL-ի վերիֆիկացիա։
  \item Vivado սինթեզ, implementation և timing analysis։
  \item LUT/FF/DSP/BRAM օգտագործման վերլուծություն։
  \item Bitstream-ի գեներացում և ամբողջական նախագծի վալիդացիա AUP ZU3-ի վրա։
  \item FPGA/PS տվյալների ինտերֆեյսի և կառուցվածքային ելքային ֆորմատի սահմանում։
\end{enumerate}

Առաջին ուսանողի դիպլոմը պետք է հիմնականում կենտրոնանա թվային նախագծման, RTL-ի, FPGA
ճարտարապետության, data/control ուղիների, հիշողության/բուֆերացման, AXI-ի,
վերիֆիկացիայի, սինթեզի, implementation-ի, timing-ի և ապարատային վալիդացիայի վրա։
